% Г-образный LC-фильтр: выпрямитель (вход слева) — дроссель — конденсатор — Rн
\documentclass[border=4pt]{standalone}
\usepackage{amsmath}
\usepackage{fontspec}
\setmainfont{Calibri}
\usepackage[european, RPvoltages]{circuitikz}
\definecolor{ink}{RGB}{31,42,51}
\definecolor{grp}{RGB}{46,110,140}
\ctikzset{bipoles/length=1.0cm}
\ctikzset{bipoles/thickness=0.9}
\ctikzset{diodes/fill=ink}
\ctikzset{inductor=american}

\begin{document}
\begin{tikzpicture}[line width=1.0pt, color=ink, font=\small]
  % --- входные зажимы выпрямителя: постоянная ЭДС E0 + пульсации u~
  \draw (0,1.4) to[sV, l_=$u_{\sim}$] (0,0)
        to[battery1, l_=$E_0$] (0,-1.4);
  % --- дроссель (переменная составляющая падает на нём)
  \draw (0,1.4) -- (0.8,1.4) to[L, l=$L_{\text{ф}}$] (2.8,1.4) node[circ]{} -- (4.4,1.4);
  % --- конденсатор фильтра
  \draw (0,-1.4) -- (2.8,-1.4);
  \draw (2.8,1.4) to[C, l_=$C_{\text{ф}}$, i=$i_C$] (2.8,-1.4);
  \draw (2.8,-1.4) -- (4.4,-1.4);
  % --- нагрузка
  \draw (4.4,1.4) to[R, l=$R_{\text{н}}$, i=$i_{\text{н}}$] (4.4,-1.4);
  % --- токи и подписи
  \node[anchor=south] at (0.35,1.5) {$i_{\text{d}}$};
  \node[anchor=south, text=grp, font=\footnotesize] at (2.2,2.0)
       {Г-образный LC-фильтр};
\end{tikzpicture}
\end{document}
